\topic{本章实验、故意破坏与练习}

\begin{experimentbox}{在模式切换两侧记录状态}
在写 CR0.PG 前和 64-bit 标签各停一次，记录 CR0、CR3、CR4、EFER、CS、RIP
和 RSP，并手算当前指令和栈的页表路径。
\end{experimentbox}

\begin{faultinjectionbox}{漏映射 Kernel 入口}
候选页表只映射 Loader，不映射 Kernel 入口。预期 far jump 后取指页故障或
重启。补上目标页映射后再切 CR3。
\end{faultinjectionbox}

\begin{pitfallbox}
先验证 ELF 再复制；先建页表再开 PG；far jump 才装入 L=1 代码段。
\end{pitfallbox}

\textbf{练习}
\begin{enumerate}[leftmargin=2.2em]
  \item 64 MiB 需要多少个 2 MiB PDE？
  \item filesz=\texttt{1800}、memsz=\texttt{2000}，清零多少 byte？
  \item CR4.PAE 是哪一位？
\end{enumerate}
\begin{answerbox}
1. 32；2. \texttt{800} byte；3. bit 5。
\end{answerbox}
